% Generated draft: requires mathematical and source review.
\documentclass[11pt]{article}
\usepackage[margin=25mm]{geometry}
\usepackage{fontspec}
\setmainfont{Latin Modern Roman}
\usepackage{amsmath,amssymb,mathtools,unicode-math}
\setmathfont{Latin Modern Math}
\usepackage{enumitem,needspace,xurl,hyperref}
\setlength{\parindent}{0pt}
\setlength{\parskip}{5pt}
\setlength{\emergencystretch}{3em}
\newcommand{\N}{\mathbb{N}}
\newcommand{\Z}{\mathbb{Z}}
\newcommand{\Q}{\mathbb{Q}}
\newcommand{\R}{\mathbb{R}}
\newcommand{\C}{\mathbb{C}}
\newcommand{\F}{\mathbb{F}}
\newcommand{\A}{\mathbb{A}}
\newcommand{\PP}{\mathbb{P}}
\newcommand{\E}{\mathbb{E}}
\newcommand{\eps}{\varepsilon}
\newcommand{\dd}{\,\mathrm{d}}
\begin{document}
\section*{UnsolvedMath ID: 20002302}
\textbf{Source problem code:} catalogue entry
\subsection{Definitions and mathematical statement}
The whole mathematical source is ``Question 18 (Jarden). Various questions from Jarden's talk.'' The trailing 3 is identified by the aid as a page number. No individual question, hypothesis, or conjecture can be recovered from this record. The aid says official sources identify the talk's theme as decidability in infinite extensions and connect Jarden--Videla's \emph{Fields on the Bottom} with the workshop project, but explicitly do not identify the missing questions.

Its developed contextual result is the following necessary condition. Fix an effective presentation of \(\overline\Q\) in which rational polynomials, algebraic-number equality, and minimal polynomials can be computed, and let
\[
 \Q\subseteq M\subseteq\overline\Q
\]
be an algebraic field. Suppose the existential ring theory of \(M\) is decidable. Then that decision procedure uniformly computes, for every irreducible \(f\in\Q[X]\),
\[
 r_M(f)=\#\{a\in M:f(a)=0\}.
\]
Indeed the finite root-count tests are existential sentences asking for \(j\le\deg f\) pairwise distinct roots; rational coefficients can be cleared. No computable membership presentation of all of \(M\) is assumed by this implication.

Define its Galois core by
\[
 \operatorname{Core}_{\Q}(M)
 =\bigcap_{\sigma\in\operatorname{Aut}(\overline\Q/\Q)}\sigma(M).
\]
For \(a\in\overline\Q\), with irreducible minimal polynomial \(f_a\), the aid's test is
\[
 a\in\operatorname{Core}_{\Q}(M)
 \quad\Longleftrightarrow\quad r_M(f_a)=\deg f_a.
\]
Thus the characteristic function of the core as a subset of the fixed effective \(\overline\Q\) is computable from the existential decision procedure. In particular a nonrecursive Galois core obstructs existential decidability, even for nonnormal \(M\). This is a necessary obstruction only, not a converse decidability criterion or an asserted reconstruction of Jarden's missing questions.
\subsection{Short English statement}
No mathematical question is recoverable from the short Jarden reference. The aid supplies a contextual infinite-extension result: an existential decision procedure computes irreducible root counts and therefore membership in the field's Galois core. A noncomputable core is an obstruction, without identifying the lost workshop questions.
\subsection{Sources}
\begin{itemize}
\item Local UnsolvedMath registry, permanent ID 20002302, exact source code catalogue entry; canonical \texttt{aim-logic-notes.json}, \texttt{notes[77]}, tag \texttt{question}; AIM Workshop Problem Lists, Logic, \emph{Definability and decidability problems in number theory}, Question 18 (Jarden). Created and updated 14 August 2026.
\item Source URL \url{https://aimath.org/pastworkshops/definabilityinntproblems.pdf}. The supplied record has no formulation beyond the talk reference; the aid identifies 3 as a page artifact. The PDF was not accessed.
\item Jarden--Videla, \emph{Fields on the Bottom}, is named by the aid as connected to the workshop's infinite-extension project. Its official-source verification, full publication entry, and identifying talk materials are not supplied. No individual problem is recovered from that thematic connection.
\item Attempt 1: \texttt{context\_only}, problem status at run \texttt{not\_a\_problem}, rigor label \texttt{proved}, novelty confidence \texttt{low}; classification \texttt{OPEN-TRIAGE}, suggested difficulty L2, registry status \texttt{open}. Candidate: root-profile and Galois-core obstruction. No separate authoritative theorem reference is supplied.
\item Ulam AI / UnsolvedMath Contributors; CC BY 4.0. Machine-generated research aid, not peer reviewed; independently verify before citation.
\end{itemize}
\end{document}
